\section{Introduction}\label{sec:intro}

Benders decomposition is a standard tool for large two-stage programs: the first-stage decisions stay in a master problem, the second stage splits into independent blocks or scenarios, and dual extreme points of the blocks generate optimality cuts that refine the master~\cite{geoffrion1970elements,birge2006introduction,rahmaniani2017benders}. In the multi-cut version, every subproblem is solved at every iteration. With hundreds of scenarios, these solves dominate the running time.

Exact solves are not needed for validity. Any dual-\emph{feasible} vector of a subproblem gives a valid, possibly weaker, cut~\cite{zakeri2000inexact,yang2012tighter}. This creates a trade-off between how strong the dual information behind a cut is and how much it costs to obtain, and it suggests a role for learning: solve a few subproblems exactly, and infer dual-feasible vectors for the others from what has been observed. Three things must hold for this to work. The predictor must use information that is actually available when a cut is needed; its output must be certified dual feasible, otherwise the cut is invalid; and the certification must be cheap, otherwise nothing is saved.

We propose \emph{Learning-Augmented Benders Decomposition} (LABD). After a short warm-up in which all subproblems are solved exactly, LABD learns a representation of each subproblem from the trajectory of its optimal duals, fixes a small set of subproblems to be solved exactly, and trains a mapping that predicts the dual vector of every other subproblem from its representation and from its current dual objective. A raw prediction may be infeasible. Restoring feasibility by optimization means either solving the subproblem itself or solving a linear program that finds a nearby dual-feasible vector; both are posed over the same dual polyhedron, so either one costs about as much as the exact solve that prediction was meant to avoid. We show that a single dual-feasible reference point removes this cost: moving from the reference point toward the prediction for as long as the dual constraints allow gives a feasible vector, and the length of this move is given by an explicit formula. The same move can be chosen so that the resulting cut is never weaker than the cut of the reference point at the current first-stage solution. We exploit the structure of the problem, a dual feasible region shared by all subproblems, to obtain such a reference point at every iteration from the subproblems that are solved exactly.

\subsection{Related Work}\label{subsec:related}

Inexact Benders methods use dual-feasible, non-optimal solutions and keep convergence under suitable conditions~\cite{zakeri2000inexact}; such cuts need not come from extreme points and may be weaker~\cite{yang2012tighter}. Sampling reduces the second-stage information evaluated per iteration~\cite{infanger1992monte}, and Bertsimas et al.~\cite{bertsimas2024stochastic} solve a random sample of scenarios and construct deterministic valid cuts for the rest by aggregating the sampled duals. LABD shares the premise that useful cuts do not require an optimal dual for every subproblem at every iteration. It differs in that the dual vector of each unsolved subproblem is predicted from subproblem-specific information and the current state, and is then certified.

Learning has been used inside Benders decomposition to propose master solutions that are then evaluated by the usual Benders steps~\cite{mana2024accelerating,agyeman2026graph}, and to decide which already-generated cuts to keep, by classification~\cite{jia2021benders} or reinforcement learning~\cite{cai2026rlbd}; related work in general cutting-plane methods learns which separators to run, when to stop, or which cuts to add~\cite{li2023separators,ling2024stop,huang2022learning,paulus2022learning}. In all of these, the subproblem information behind a cut must still be computed. Closer to LABD is work that learns the cut-generation oracle itself. Neural Benders learns which optimal dual to use for a degenerate subproblem that is being solved~\cite{monemi2025neuralbenders}. Learned value functions and multipliers yield approximate generalized Benders cuts~\cite{mitrai2024mlbenders}; KKT-informed networks predict primal--dual subproblem solutions and use them as inexact cuts~\cite{agyeman2026hybrid}; and learned duals with a safe shift have been used in nested Benders~\cite{aslani2025lcc}. Learned recourse functions have also been embedded directly in two-stage models~\cite{patel2022neur2sp}.  Larsen et al.~\cite{larsen2024fast} replace exact second-stage solutions in the L-shaped method by supervised predictions trained in advance on similar instances, which yields fast heuristics with a small loss in accuracy; LABD instead learns within the instance being solved, from its own warm-up, keeps a subset of subproblems exact, and certifies every predicted dual so that all cuts remain valid.

When learned information defines a cut, validity becomes the central issue. Proxy-BD~\cite{guan2026proxy} learns a stand-alone proxy of the subproblem response and passes its output through a predict--project--complete layer that returns dual-feasible certificates. LABD also separates validity from prediction accuracy, but its certificate needs no projection: feasibility is restored by a single explicit step toward the prediction from a reference point built from the exact duals of the current iteration, a step that exploits the relationships among subproblems of the same instance. The certificate also comes with a bound on the resulting loss in solution quality.

\subsection{Our Contributions}\label{subsec:contributions}
\begin{itemize}
\item We extract the learning information from the Benders process itself: during warm-up, exact dual responses are collected across subproblems and first-stage solutions, and their trajectories yield subproblem representations through graph embeddings, with little reliance on externally supplied features.
\item We develop LABD, which solves a representative subset of subproblems exactly and generates dual-feasible predictions for the others within one Benders procedure, so that every generated cut remains valid.
\item We give a certification step that turns any prediction into a dual-feasible vector without solving an optimization problem, using a reference point formed from the current exact duals, and chooses the step so that the predicted cut is never weaker than the cut of the reference point where it is generated.
\item We establish optimality-gap bounds that separate a structural term, which depends only on how far the predicted subproblems are from the exact ones, from a nonnegative term contributed by learning, together with a gap certificate that can be computed at the returned solution.
\item We evaluate LABD on stochastic capacitated facility location instances with 100 to 500 scenarios. Compared with Standard BD, LABD reduces the median solution time by about 78\% (a speedup of about $4.6\times$) and returns the optimal objective in 34 of the 36 configurations that Standard BD solves, and it is faster than the random-selection benchmark that solves the same number of subproblems exactly in 33 of them.
\end{itemize}

